
%%%%%%%%%%%%%%%%%%%%%%%%%%%%%%%%%%%%%%%%%%%%%%%%%%%%%%%%%%%%%%%%%%%%%%%
%                   ANFANG ABSTRACT
%%%%%%%%%%%%%%%%%%%%%%%%%%%%%%%%%%%%%%%%%%%%%%%%%%%%%%%%%%%%%%%%%%%%%%%
%
% Autor: Felix Zauner, Franz Spatzenegger
%
%
% WICHTIG:
%   Bitte füllen Sie die imfolgenden mit "AUSFÜLLEN" gekennzeichneten Felder gewissenhaft aus.
%   Felder mit "NICHT AUSFÜLLEN" sind vom Benutzer nicht zu bearbeiten da schon vorausgefüllt.
%
% Der in den Feldern stehende Text dient nur als Beispiel und muss entsprechend angepasst werden.
%
%
% Erneuter Hinweis bezüglich Einbindung der benötigten Pakete für Abstract
%
%   Benötigte Pakete für Abstract
%       \usepackage{graphicx}       % für Grafiken
%       \usepackage[T1]{fontenc}    % für Umlaute und Sonderzeichen
%
% 
% Bei Problemen jeglicher Art wenden Sie sich bitte an einen der Autoren dieser Vorlage.
%
%
%%%%%%%%%%%%%%%%%%%%%%%%%%%%%%%%%%%%%%%%%%%%%%%%%%%%%%%%%%%%%%%%%%%%%%%   


\begin{center}
    {\fontsize{30}{36}\selectfont\bfseries Diplomarbeit}\\[0.5em]
    {\fontsize{20.74}{24}\selectfont\bfseries Dokumentation}
\end{center}

\vspace{0.5cm}
\noindent
\begin{tabular}{|p{0.31\textwidth}|p{0.61\textwidth}|}
\hline
% AUSFÜLLEN: Tragen Sie hier alle Autorennamen ein. Format: Vorname Nachname, ggf. mehrere durch Komma getrennt.
Name der Verfasser & Felix Zauner, Timofey Luzin \\
\hline
% NICHT AUSFÜLLEN: Jahrgang/Schulform (z. B. 5AHET). Wenn nicht zutreffend, anpassen oder weglassen.
Jahrgang & 5AHETS \\
% NICHT AUSFÜLLEN: Schuljahr im Format JJJJ/JJ (z. B. 2025/26).
Schuljahr & 2025/26  \\
\hline
% AUSFÜLLEN: Prägnanter Titel des Themas. Sollte die Kernidee der Arbeit widerspiegeln.
Thema der Diplomarbeit & ALBERT - Entwicklung eines autonomen Lenkflugkörpers mit Echtzeit-Fluglagenregelung zur Wetterdatenerfassung \\
\hline
% AUSFÜLLEN: Externe Partner oder Firmen, die an der Arbeit beteiligt sind. Falls keine, "keine" oder leer lassen.
Kooperationspartner & Palfinger AG, Extrudr GmbH\\
\hline
\noalign{\vskip 1mm}
\hline
% AUSFÜLLEN: Hier die Aufgabenstellung in ein paar Sätzen beschreiben.
Aufgabenstellung & Ziel der Arbeit ist die Entwicklung einer flugfähigen, finnengesteuerten Rakete mit aktiver Flugregelung sowie eines eigens entwickelten Flightcomputers, der eine autonome Flugführung zu vorgegebenen GPS-Zielpunkten ermöglicht. \\
\hline
\noalign{\vskip 1mm}
\hline
% AUSFÜLLEN: Beschreibung der Realisierung. Hier sollte die Umsetzung des Projektes beschrieben werden.
Realisierung & Zur Bestimmung der Fluglage werden auf der Platine Sensoren implementiert und deren Daten mittels Senorfusion verarbeitet.
Auf Basis einer Zustandsraummodellierung der Rakete wird ein LQR-Regler entworfen, der die erforderlichen Stellgrößen zur Ansteuerung der Finnen berechnet und so eine Stabilisierung sowie gezielte Flugbahnführung ermöglicht. \\
\hline
\noalign{\vskip 1mm}
\hline
% AUSFÜLLEN: Konkrete Ergebnisse zusammenfassen
Ergebnisse & Im Rahmen der Arbeit wurde ein Gesamtsystem bestehend aus Finnensteuerung, Fallschirmauswurf und Flightcomputer entwickelt. 
Die Regelung wurde zunächst in Simulationen getestet und anschließend auf dem Flightcomputer implementiert, wobei 
Hardware-in-the-Loop-Tests die grundlegende Funktionalität bestätigten. \\
\hline
\noalign{\vskip 1mm}
\hline
% AUSFÜLLEN: Fügen Sie hier die Bilddatei und eine kurze Bildunterschrift ein. Bilddateien sollten im gleichen Verzeichnis liegen oder mit Pfad angegeben werden.
Foto (mit Eläuterung) & \vspace{0.1cm}{\centering\includegraphics[width=0.5\linewidth]{ski.jpg}\par} Der alpenländische Warp-Antrieb in Aktion.\\
\hline
\noalign{\vskip 1mm}
\hline
% NICHT AUSFÜLLEN: Geben Sie an, wo die Arbeit eingesehen werden kann (physisch und digital). Nennen Sie Ansprechpartner und ggf. Kontaktinformationen.
Möglichkeit der Einsichtnahme in die Arbeit & Die Diplomarbeit ist in gebundener Form sowohl in der Schulbibliothek als auch bei AV Prof. Dipl-Ing. (FH) Roland Holzer einzusehen. Darüber hinaus besitzt jedes Mitglied des Projektteams eine vollständige Version in gebundener und digitaler Form.\\
\hline
\noalign{\vskip 1mm}
\hline
% NICHT AUSFÜLLEN: Unterschriftenfelder: Datum/Unterschrift, Prüfer/in, Abteilungsvorstand. Diese Felder bleiben leer und werden manuell unterschrieben.
\multicolumn{2}{|c|}{%
    \begin{tabular}{p{0.30\textwidth}|p{0.30\textwidth}|p{0.30\textwidth}}
        Approbation\newline(Datum/Unterschrift)\newline\newline & Prüfer/Prüferin & Abteilungsvorstand
    \end{tabular}
} \\
\hline
\end{tabular}

\newpage

\begin{center}
    {\fontsize{30}{36}\selectfont\bfseries Diploma Thesis}\\[0.5em]
    {\fontsize{20.74}{24}\selectfont\bfseries Documentation}
\end{center}
\vspace{0.5cm}
\noindent
\begin{tabular}{|p{0.31\textwidth}|p{0.61\textwidth}|}
\hline
% AUSFÜLLEN: Tragen Sie hier alle Autorennamen ein. Format: Vorname Nachname, ggf. mehrere durch Komma getrennt.
Author(s) & Felix Zauner, Timofey Luzin \\
\hline
% NICHT AUSFÜLLEN: Jahrgang/Schulform (z. B. 5AHET). Wenn nicht zutreffend, anpassen oder weglassen.
Form & 5AHET \\
% NICHT AUSFÜLLEN: Schuljahr im Format JJJJ/JJ (z. B. 2025/26).
Academic year & 2025/26  \\
\hline
% AUSFÜLLEN: Prägnanter Titel des Themas IN ENGLISCH. Sollte die Kernidee der Arbeit widerspiegeln.
Topic & ALBERT - Development of an Autonomous Guided Missile with Real-Time Attitude Control for Weather Data Collection \\
\hline
% AUSFÜLLEN: Externe Partner oder Firmen, die an der Arbeit beteiligt sind. Falls keine, "keine" oder leer lassen.
Cooperation Partners & Palfinger AG, Extrudr GmbH \\
\hline
\noalign{\vskip 1mm}
\hline
% AUSFÜLLEN: Hier die Aufgabenstellung IN ENGLISCH in ein paar Sätzen beschreiben.
Assignment & The assignment is the development of a prototype for a warp drive using modern propulsion technologies \\
\hline
\noalign{\vskip 1mm}
\hline
% AUSFÜLLEN: Beschreibung der Realisierung IN ENGLISCH. Hier sollte die Umsetzung des Projektes beschrieben werden.
Realization & A warp drive born of Alpine engineering relies less on physics and more on intuition, charm, and improvisation. Spacetime resembles a bumpy ski run, its energy supplied by a crate of beer and a great deal of persuasion. Calculations are produced as rhymed Gstanzln (folk verses) rather than in the laboratory. Thus the universe moves -- not by technology, but by creativity and the belief that even the cosmos can be softened with a smile.\\
\hline
\noalign{\vskip 1mm}
\hline
% AUSFÜLLEN: Konkrete Ergebnisse IN ENGLISCH zusammenfassen
Result & The developed warp drive enables interstellar journeys (Stubaital) to be completed in record time, with Alpine engineering providing a unique blend of efficiency and charm. Early tests indicate that the drive is not only functional but also has the potential to revolutionize the space industry. \\
\hline
\noalign{\vskip 1mm}
\hline
% AUSFÜLLEN: Fügen Sie hier die Bilddatei und eine kurze Bildunterschrift ein. Bilddateien sollten im gleichen Verzeichnis liegen oder mit Pfad angegeben werden.
Illustrative Graph, Photo (with explanation) & \vspace{0.1cm}{\centering\includegraphics[width=0.5\linewidth]{ski.jpg}\par} The Alpine warp drive in action.\\
\hline
\noalign{\vskip 1mm}
\hline
% NICHT AUSFÜLLEN: Geben Sie an, wo die Arbeit eingesehen werden kann (physisch und digital). Nennen Sie Ansprechpartner und ggf. Kontaktinformationen.
Accessibility of \newline Diploma Thesis & The diploma thesis is available in the school library as well as in the office of the head of department Prof. Dipl-Ing. (FH) Roland Holzer. In addition, each member of the project team has a complete version in hardcover and digital form.\\
\hline
\noalign{\vskip 1mm}
\hline
% NICHT AUSFÜLLEN: Unterschriftenfelder: Datum/Unterschrift, Prüfer/in, Abteilungsvorstand. Diese Felder bleiben leer und werden manuell unterschrieben.
\multicolumn{2}{|l|}{%
    \begin{tabular}{p{0.30\textwidth}|p{0.30\textwidth}|p{0.30\textwidth}}
        Approval\newline(Date/Signature)\newline\newline &  Examiner & Head of Department
    \end{tabular}
} \\
\hline
\end{tabular}
